% Chapter 1, Figure 10: aerodynamic (a-c) and mechanical (d-f) disturbances (scan: figures/ch1/fig10.png,
% PDF 79). Six copies of the model rocket, nose up, C.G. above C.P.; the disturbing force as a heavy arrow.
\documentclass[11pt]{standalone}
\usepackage{tamrfig}
\begin{document}
\begin{tikzpicture}
  \def\L{3.6}\def\r{0.17}\def\cr{0.75}\def\ct{0.25}\def\sp{0.55}\def\sw{0.5}
  \tikzset{push/.style={draw=ink, line width=2.2pt, -{Stealth[length=7pt,width=6pt]}},
           wind/.style={draw=ink2, line width=0.6pt, -{Stealth[length=4pt,width=2.8pt]}}}
  % \dist{<x>}{<y>}{<plume>}{<extra rocket keys>}: a rocket with its tail at (x,y); C.G. and C.P. marked
  \newcommand{\dist}[4]{%
    \pic[rotate=-90] at (#1,#2+\L) {rocket={model, length=\L, diameter={2*\r}, nose length=0.9,
        root chord=\cr, tip chord=\ct, span=\sp, sweep=\sw, #4}};
    \ifdim #3pt>0pt \draw[thin line, shift={(#1,#2)}, rotate=-90] \plumeshapeb{1.25}{0.6*\r}{#3}; \fi
    \node[cg mark] at (#1,#2+0.36*\L) {};
    \node[cp mark] at (#1,#2+0.22*\L) {};}
  % (a) horizontal wind, acting at the C.P.
  \dist{0}{0}{1.1}{}
  \foreach \y in {-0.75,-0.2,0.35,1.45,2.0,2.55,3.1,3.65} \draw[wind] (1.45,\y) -- (1.05,\y);
  \draw[push] (1.25,0.22*\L) -- (0.2,0.22*\L);
  \node[panel] at (1.4,-1.25) {(a)};
  % (b) a mismounted fin: its root canted across the body
  \begin{scope}[shift={(3.4,0)}]
    \dist{0}{0}{1.1}{}
    \draw[outline] (-0.04,0.66) -- (0.07,-0.02);
    \draw[push] (1.25,0.35) -- (0.45,0.35);
    \node[panel] at (1.4,-1.25) {(b)};
  \end{scope}
  % (c) torsionally vibrating fins: their deflected positions, pushed from both sides
  \begin{scope}[shift={(6.8,0)}]
    \dist{0}{0}{1.1}{}
    \foreach \s in {1,-1} \draw[thin line, dashed] (\s*\r,0.75) -- (\s*0.62,0.15) -- (\s*0.72,0) (\s*\r,0.75) -- (\s*0.8,0.3) -- (\s*0.7,0);
    \draw[thin line] (-0.1,0.7) -- (0.1,0.05) (0.1,0.7) -- (-0.1,0.05);
    \draw[push] (-1.35,0.35) -- (-0.6,0.35);
    \draw[push] (1.35,0.35) -- (0.6,0.35);
    \node[panel] at (1.4,-1.25) {(c)};
  \end{scope}
  % (d) a mismounted engine: the plume and the thrust line askew
  \begin{scope}[shift={(0,-6.3)}]
    \dist{0}{0}{0}{}
    \draw[thin line, rotate=8, rotate=-90] \plumeshapeb{1.25}{0.6*\r}{1.2};
    \draw[thin line] (0,0) -- (0,-1.3);
    \draw[thin line] (-0.28,2.4) -- (0.2,-1.3);
    \draw[push] (1.25,0.02) -- (0.2,0.02);
    \node[panel] at (1.4,-1.25) {(d)};
  \end{scope}
  % (e) launch rod whip: the lug leaves the rod
  \begin{scope}[shift={(3.4,-6.3)}]
    \dist{0}{0}{1.1}{}
    \draw[outline] (\r,0.36*\L+0.05) rectangle (\r+0.07,0.36*\L+0.35);
    \draw[outline] (\r+0.035,0.45) -- (\r+0.035,-1.3);
    \draw[thin line] (\r-0.065,-1.3) -- (\r+0.135,-1.3);
    \draw[push] (1.25,0.02) -- (0.3,0.02);
    \node[panel] at (1.4,-1.25) {(e)};
  \end{scope}
  % (f) staging: the lower stage blown to the right, the upper stage's base pushed to the left
  \begin{scope}[shift={(6.8,-6.3)}]
    \dist{0}{0}{0}{}
    \begin{scope}[shift={(0.2,-0.18)}, rotate=-20]
      \draw[outline, fill=white] (-\r,0) rectangle (\r,-0.85);
      \foreach \s in {1,-1} \draw[outline] (\s*\r,-0.15) -- (\s*0.85,-0.45) -- (\s*0.85,-0.85) -- (\s*\r,-0.85);
      \draw[centerline] (0,0.2) -- (0,-1.2);
    \end{scope}
    \foreach \a in {0,20,...,340} \draw[thin line, draw=ink2] (\a:0.12) -- (\a:{0.35+0.12*mod(\a,40)/20});
    \draw[push] (1.35,0.05) -- (0.35,0.05);
    \node[panel] at (1.4,-1.25) {(f)};
  \end{scope}
\end{tikzpicture}
\end{document}
